% Options for packages loaded elsewhere
\PassOptionsToPackage{unicode}{hyperref}
\PassOptionsToPackage{hyphens}{url}
%
\documentclass[
]{article}
\usepackage{amsmath,amssymb}
\usepackage{iftex}
\ifPDFTeX
  \usepackage[T1]{fontenc}
  \usepackage[utf8]{inputenc}
  \usepackage{textcomp} % provide euro and other symbols
\else % if luatex or xetex
  \usepackage{unicode-math} % this also loads fontspec
  \defaultfontfeatures{Scale=MatchLowercase}
  \defaultfontfeatures[\rmfamily]{Ligatures=TeX,Scale=1}
\fi
\usepackage{lmodern}
\ifPDFTeX\else
  % xetex/luatex font selection
\fi
% Use upquote if available, for straight quotes in verbatim environments
\IfFileExists{upquote.sty}{\usepackage{upquote}}{}
\IfFileExists{microtype.sty}{% use microtype if available
  \usepackage[]{microtype}
  \UseMicrotypeSet[protrusion]{basicmath} % disable protrusion for tt fonts
}{}
\makeatletter
\@ifundefined{KOMAClassName}{% if non-KOMA class
  \IfFileExists{parskip.sty}{%
    \usepackage{parskip}
  }{% else
    \setlength{\parindent}{0pt}
    \setlength{\parskip}{6pt plus 2pt minus 1pt}}
}{% if KOMA class
  \KOMAoptions{parskip=half}}
\makeatother
\usepackage{xcolor}
\usepackage[margin=1in]{geometry}
\usepackage{color}
\usepackage{fancyvrb}
\newcommand{\VerbBar}{|}
\newcommand{\VERB}{\Verb[commandchars=\\\{\}]}
\DefineVerbatimEnvironment{Highlighting}{Verbatim}{commandchars=\\\{\}}
% Add ',fontsize=\small' for more characters per line
\usepackage{framed}
\definecolor{shadecolor}{RGB}{248,248,248}
\newenvironment{Shaded}{\begin{snugshade}}{\end{snugshade}}
\newcommand{\AlertTok}[1]{\textcolor[rgb]{0.94,0.16,0.16}{#1}}
\newcommand{\AnnotationTok}[1]{\textcolor[rgb]{0.56,0.35,0.01}{\textbf{\textit{#1}}}}
\newcommand{\AttributeTok}[1]{\textcolor[rgb]{0.13,0.29,0.53}{#1}}
\newcommand{\BaseNTok}[1]{\textcolor[rgb]{0.00,0.00,0.81}{#1}}
\newcommand{\BuiltInTok}[1]{#1}
\newcommand{\CharTok}[1]{\textcolor[rgb]{0.31,0.60,0.02}{#1}}
\newcommand{\CommentTok}[1]{\textcolor[rgb]{0.56,0.35,0.01}{\textit{#1}}}
\newcommand{\CommentVarTok}[1]{\textcolor[rgb]{0.56,0.35,0.01}{\textbf{\textit{#1}}}}
\newcommand{\ConstantTok}[1]{\textcolor[rgb]{0.56,0.35,0.01}{#1}}
\newcommand{\ControlFlowTok}[1]{\textcolor[rgb]{0.13,0.29,0.53}{\textbf{#1}}}
\newcommand{\DataTypeTok}[1]{\textcolor[rgb]{0.13,0.29,0.53}{#1}}
\newcommand{\DecValTok}[1]{\textcolor[rgb]{0.00,0.00,0.81}{#1}}
\newcommand{\DocumentationTok}[1]{\textcolor[rgb]{0.56,0.35,0.01}{\textbf{\textit{#1}}}}
\newcommand{\ErrorTok}[1]{\textcolor[rgb]{0.64,0.00,0.00}{\textbf{#1}}}
\newcommand{\ExtensionTok}[1]{#1}
\newcommand{\FloatTok}[1]{\textcolor[rgb]{0.00,0.00,0.81}{#1}}
\newcommand{\FunctionTok}[1]{\textcolor[rgb]{0.13,0.29,0.53}{\textbf{#1}}}
\newcommand{\ImportTok}[1]{#1}
\newcommand{\InformationTok}[1]{\textcolor[rgb]{0.56,0.35,0.01}{\textbf{\textit{#1}}}}
\newcommand{\KeywordTok}[1]{\textcolor[rgb]{0.13,0.29,0.53}{\textbf{#1}}}
\newcommand{\NormalTok}[1]{#1}
\newcommand{\OperatorTok}[1]{\textcolor[rgb]{0.81,0.36,0.00}{\textbf{#1}}}
\newcommand{\OtherTok}[1]{\textcolor[rgb]{0.56,0.35,0.01}{#1}}
\newcommand{\PreprocessorTok}[1]{\textcolor[rgb]{0.56,0.35,0.01}{\textit{#1}}}
\newcommand{\RegionMarkerTok}[1]{#1}
\newcommand{\SpecialCharTok}[1]{\textcolor[rgb]{0.81,0.36,0.00}{\textbf{#1}}}
\newcommand{\SpecialStringTok}[1]{\textcolor[rgb]{0.31,0.60,0.02}{#1}}
\newcommand{\StringTok}[1]{\textcolor[rgb]{0.31,0.60,0.02}{#1}}
\newcommand{\VariableTok}[1]{\textcolor[rgb]{0.00,0.00,0.00}{#1}}
\newcommand{\VerbatimStringTok}[1]{\textcolor[rgb]{0.31,0.60,0.02}{#1}}
\newcommand{\WarningTok}[1]{\textcolor[rgb]{0.56,0.35,0.01}{\textbf{\textit{#1}}}}
\usepackage{graphicx}
\makeatletter
\newsavebox\pandoc@box
\newcommand*\pandocbounded[1]{% scales image to fit in text height/width
  \sbox\pandoc@box{#1}%
  \Gscale@div\@tempa{\textheight}{\dimexpr\ht\pandoc@box+\dp\pandoc@box\relax}%
  \Gscale@div\@tempb{\linewidth}{\wd\pandoc@box}%
  \ifdim\@tempb\p@<\@tempa\p@\let\@tempa\@tempb\fi% select the smaller of both
  \ifdim\@tempa\p@<\p@\scalebox{\@tempa}{\usebox\pandoc@box}%
  \else\usebox{\pandoc@box}%
  \fi%
}
% Set default figure placement to htbp
\def\fps@figure{htbp}
\makeatother
\setlength{\emergencystretch}{3em} % prevent overfull lines
\providecommand{\tightlist}{%
  \setlength{\itemsep}{0pt}\setlength{\parskip}{0pt}}
\setcounter{secnumdepth}{-\maxdimen} % remove section numbering
\usepackage{bookmark}
\IfFileExists{xurl.sty}{\usepackage{xurl}}{} % add URL line breaks if available
\urlstyle{same}
\hypersetup{
  pdftitle={Lab 7 - Inference on Means},
  hidelinks,
  pdfcreator={LaTeX via pandoc}}

\title{Lab 7 - Inference on Means}
\author{}
\date{\vspace{-2.5em}}

\begin{document}
\maketitle

\subsection{North Carolina Births}\label{north-carolina-births}

In 2004, the state of North Carolina released to the public a large data
set containing information on births recorded in this state. This data
set has been of interest to medical researchers who are studying the
relation between habits and practices of expectant mothers and the birth
of their children. We will work with a sample of observations from this
data set. These 999 cases were chosen at random.

Let's load the \texttt{nc} data set into our workspace. The data is in a
file called \texttt{nc.csv}.

\vspace{3mm}

\begin{Shaded}
\begin{Highlighting}[]
\CommentTok{\# change working directory to the folder}
\CommentTok{\# where the data is saved}
\NormalTok{nc }\OtherTok{=} \FunctionTok{read.csv}\NormalTok{(}\StringTok{"nc.csv"}\NormalTok{, }\AttributeTok{header=}\NormalTok{T)}
\CommentTok{\# attach dataset to workspace}
\FunctionTok{attach}\NormalTok{(nc)}
\end{Highlighting}
\end{Shaded}

We have observations on 13 different variables, some categorical and
some numerical. The meaning of each variable is as follows.

\begin{table}[h] \small
\begin{tabular}{r | l}
\texttt{fage} & father's age in years. \\
\texttt{mage} & mother's age in years. \\
\texttt{mature} & maturity status of mother. \\
\texttt{weeks} & length of pregnancy in weeks. \\
\texttt{premie} & whether the birth was classified as premature (premie) or full-term. \\
\texttt{visits} & number of hospital visits during pregnancy. \\
\texttt{gained} & weight gained by mother during pregnancy in pounds. \\
\texttt{weight} & weight of the baby at birth in pounds. \\
\texttt{lowbirthweight} & whether baby was classified as low birthweight (`low`) or not (`not low`). \\
\texttt{gender} & gender of the baby, `female` or `male`. \\
\texttt{habit} & status of the mother as a `nonsmoker` or a `smoker`. \\
\texttt{marital} & whether mother is `married` or `not married` at birth. \\
\texttt{whitemom} & whether mom is `white` or `not white`. \\
\end{tabular}
\end{table}

First, let's calculate the mean and the standard deviation of babies'
weights in this sample as well as the sample size.

\begin{Shaded}
\begin{Highlighting}[]
\NormalTok{xbar }\OtherTok{\textless{}{-}} \FunctionTok{mean}\NormalTok{(weight)}
\NormalTok{s }\OtherTok{\textless{}{-}} \FunctionTok{sd}\NormalTok{(weight)}
\NormalTok{n }\OtherTok{\textless{}{-}} \FunctionTok{length}\NormalTok{(weight)}
\end{Highlighting}
\end{Shaded}

\textbf{Exercise 1:} Make a histogram of weight of the baby,
\texttt{weight} variable, and describe the distribution of this
variable.

Now we make a density histogram to use as the backdrop and use the
\texttt{lines()} function to overlay a normal probability curve. The
difference between a frequency histogram and a density histogram is that
while in a frequency histogram the \texttt{heights} of the bars add up
to the total number of observations, in a density histogram the
\texttt{areas} of the bars add up to 1. The area of each bar can be
calculated as simply the height \(\times\) the width of the bar. Using a
density histogram allows us to properly overlay a normal distribution
curve over the histogram since the curve is a normal probability density
function. Frequency and density histograms both display the same exact
shape; they only differ in their y-axis. You can verify this by
comparing the frequency histogram you constructed earlier and the
density histogram created by the commands below.

\begin{Shaded}
\begin{Highlighting}[]
\CommentTok{\# Plot histogram}
\FunctionTok{hist}\NormalTok{(weight, }\AttributeTok{probability =} \ConstantTok{TRUE}\NormalTok{)}
\CommentTok{\# Weights range }
\NormalTok{x }\OtherTok{=} \FunctionTok{seq}\NormalTok{(}\DecValTok{1}\NormalTok{, }\DecValTok{13}\NormalTok{, }\AttributeTok{by=}\FloatTok{0.1}\NormalTok{)}
\CommentTok{\# Compute density for x values}
\NormalTok{y }\OtherTok{=} \FunctionTok{dnorm}\NormalTok{(}\AttributeTok{x =}\NormalTok{ x, }\AttributeTok{mean =}\NormalTok{ xbar, }\AttributeTok{sd =}\NormalTok{ s)}
\CommentTok{\# Draw line on top of histogram}
\FunctionTok{lines}\NormalTok{(}\AttributeTok{x =}\NormalTok{ x, }\AttributeTok{y =}\NormalTok{ y, }\AttributeTok{col =} \StringTok{"blue"}\NormalTok{)}
\end{Highlighting}
\end{Shaded}

\begin{center}\includegraphics[width=0.47\linewidth]{Lab7-HW5_files/figure-latex/unnamed-chunk-3-1} \end{center}

After plotting the density histogram with the first command, we create
the x- and y-coordinates for the normal curve. We chose the \texttt{x}
range as 1 to 13 in order to span the entire range of \texttt{weight}.
To create \texttt{y}, we use \texttt{dnorm()} to calculate the density
of each of those x-values in a distribution that is normal with mean
\texttt{xbar} and standard deviation \texttt{s}. The final command draws
a curve on the existing plot (the density histogram) by connecting each
of the points specified by \texttt{x} and \texttt{y}.

Eyeballing the shape of the histogram is one way to determine if the
data appear to be nearly normally distributed. A normal probability plot
or Q-Q plot is also a good alternative approach to check for normality.

\begin{Shaded}
\begin{Highlighting}[]
\FunctionTok{qqnorm}\NormalTok{(weight)}
\FunctionTok{qqline}\NormalTok{(weight)}
\end{Highlighting}
\end{Shaded}

\begin{center}\includegraphics[width=0.45\linewidth]{Lab7-HW5_files/figure-latex/unnamed-chunk-4-1} \end{center}

\textbf{Exercise 2:} Based on these plots, does it appear that the data
follow a nearly normal distribution?

Note that event though the Q-Q plot displays a deviation from normality,
the sample size is large enough (\(n=999 >30\)) such that the results
from the Central Limit Theorem still apply and we can proceed to perform
inference on the parameters of interest.

\subsection{Confidence Intervals for a
Mean}\label{confidence-intervals-for-a-mean}

If we want to estimate the average weight of a baby at birth, we can use
information from this sample to build a confidence interval for the true
population parameter.

Remember that a confidence interval is always of the form

\begin{center}
$ point~estimate + ME$
\end{center}

In this case the point estimate is the sample mean. The margin of error
is always composed of two components, a critical value (\(t^\star\) or
\(z^\star\)) and a standard error multiplied together. The standard
error for the sample mean can be calculated as \(\frac{s}{\sqrt{n}}\).
Let's calculate that first.

\begin{Shaded}
\begin{Highlighting}[]
\NormalTok{se }\OtherTok{\textless{}{-}}\NormalTok{ s }\SpecialCharTok{/} \FunctionTok{sqrt}\NormalTok{(n)}
\end{Highlighting}
\end{Shaded}

Note that for \(n=999\), the t-distribution -- which is the exact
distribution to use to calculate a confidence interval for the mean --
is indistinguishable from the normal. By now, you probably know that
\(z^\star\) for a 95\% confidence interval is 1.96. This value comes
from the fact that \(z = -1.96\) and \(z = 1.96\) are the cutoff values
for the middle 95\% of a standard normal distribution. In other words,
they are the cutoff points for the 2.5\(^{th}\) percentile (in the lower
end) and \(2.5 + 95 = 97.5^{th}\) percentile (in the upper end). In
order to get R to calculate this value for us, we use

\begin{Shaded}
\begin{Highlighting}[]
\FunctionTok{qnorm}\NormalTok{(}\FloatTok{0.975}\NormalTok{)}
\end{Highlighting}
\end{Shaded}

\begin{verbatim}
## [1] 1.959964
\end{verbatim}

So if we wanted to calculate the critical value for a 98\% confidence
interval, we first need to realize that the middle 98\% of the
distribution will be bounded by the 1\(^{st}\) and the 99\(^{th}\)
percentiles, and hence the critical value will be

\begin{Shaded}
\begin{Highlighting}[]
\FunctionTok{qnorm}\NormalTok{(}\FloatTok{0.99}\NormalTok{)}
\end{Highlighting}
\end{Shaded}

\begin{verbatim}
## [1] 2.326348
\end{verbatim}

since the critical value is always positive.

\textbf{Exercise 3:} Verify that the critical values above are very
similar to those obtained from the corresponding t-distribution.

Let's get back to the calculation of the confidence interval. If we want
to construct a 95\% confidence interval we can set \(z^\star\) to be

\begin{Shaded}
\begin{Highlighting}[]
\NormalTok{zstar }\OtherTok{\textless{}{-}} \FunctionTok{qnorm}\NormalTok{(}\FloatTok{0.975}\NormalTok{)}
\end{Highlighting}
\end{Shaded}

The margin of error is simply the product of \(z^\star\) and the
standard error

\begin{Shaded}
\begin{Highlighting}[]
\NormalTok{me }\OtherTok{\textless{}{-}}\NormalTok{ zstar }\SpecialCharTok{*}\NormalTok{ se}
\end{Highlighting}
\end{Shaded}

Then, the bounds of the confidence interval can be calculated as

\begin{Shaded}
\begin{Highlighting}[]
\NormalTok{xbar }\SpecialCharTok{{-}}\NormalTok{ me}
\end{Highlighting}
\end{Shaded}

\begin{verbatim}
## [1] 7.011111
\end{verbatim}

\begin{Shaded}
\begin{Highlighting}[]
\NormalTok{xbar }\SpecialCharTok{+}\NormalTok{ me}
\end{Highlighting}
\end{Shaded}

\begin{verbatim}
## [1] 7.197838
\end{verbatim}

\textbf{Exercise 4:} Interpret the confidence interval. Make sure to put
it in context of the data.

\subsection{Hypothesis Testing for
Means}\label{hypothesis-testing-for-means}

In first world nations, the average birth weight of a full-term new born
is approximately 7.5 pounds. Let's see how the average weight of babies
born in North Carolina compare to this value.

In order to do this we need to first subset our dataset to include only
full-term babies. We'll call this new data set \texttt{nc\_ft}.

\begin{Shaded}
\begin{Highlighting}[]
\NormalTok{nc\_ft }\OtherTok{\textless{}{-}} \FunctionTok{subset}\NormalTok{(nc, nc}\SpecialCharTok{$}\NormalTok{premie }\SpecialCharTok{==} \StringTok{"full term"}\NormalTok{)}
\end{Highlighting}
\end{Shaded}

You should see a new data set appear in your workspace which has 846
rows and 13 variables. Earlier we saw that there were 152 premature and
846 full term births, hence we know that there should be 846 cases
(rows) in the new data set that we just created.

Let's calculate the mean birth weight for the full term babies

\begin{Shaded}
\begin{Highlighting}[]
\NormalTok{xbar }\OtherTok{=} \FunctionTok{mean}\NormalTok{(nc\_ft}\SpecialCharTok{$}\NormalTok{weight)}
\NormalTok{boot\_means }\OtherTok{\textless{}{-}} \FunctionTok{replicate}\NormalTok{(}\DecValTok{5000}\NormalTok{, }\FunctionTok{mean}\NormalTok{(}\FunctionTok{sample}\NormalTok{(nc\_ft}\SpecialCharTok{$}\NormalTok{weight, }\AttributeTok{replace =} \ConstantTok{TRUE}\NormalTok{)))}
\end{Highlighting}
\end{Shaded}

This value is slightly below 7.5 pounds, but it's close.

\textbf{Exercise 5:} How would you expect the sample mean of birth
weight to differ between the complete dataset (\texttt{nc}) and the
subsetted dataset (\texttt{nc\_ft})?

\textbf{Solution} The mean birth rate from the complete data set
(\texttt{nc}) should be slightly lower than the sample mean birth weight
since in the subset data set (\texttt{nc\_ft}) we excluded all of the
data points that come from babies that were not full term at the time of
birth. Since the average birth weight of a full-term new born is
approximately 7.5 pounds and the data points that we removed include all
babies that were not full term. These babies have weights that are
likely below the average weight of 7.5 pounds meaning that the sample
mean will be higher than if we did not create a subset.

\textbf{Exercise 6:} Write the hypotheses for testing if the average
weight of full term babies born in North Carolina is different than 7.5
ounces.

To conduct the hypothesis test we need to first calculate a test
statistic, a \(z\)-score, using the formula

\[ z = \frac{\bar{x} - \mu}{SE} \]

\begin{Shaded}
\begin{Highlighting}[]
\DocumentationTok{\#\# xbar = mean(nc\_ft$weight) from above \^{}\^{}}
\NormalTok{mu }\OtherTok{\textless{}{-}} \FloatTok{7.5}
\NormalTok{s }\OtherTok{\textless{}{-}} \FunctionTok{sd}\NormalTok{(nc\_ft}\SpecialCharTok{$}\NormalTok{weight)}
\NormalTok{n }\OtherTok{\textless{}{-}} \FunctionTok{length}\NormalTok{(nc\_ft}\SpecialCharTok{$}\NormalTok{weight)}
\NormalTok{se }\OtherTok{\textless{}{-}}\NormalTok{ s }\SpecialCharTok{/} \FunctionTok{sqrt}\NormalTok{(n)}
\NormalTok{z }\OtherTok{\textless{}{-}}\NormalTok{ (xbar }\SpecialCharTok{{-}}\NormalTok{ mu) }\SpecialCharTok{/}\NormalTok{ se}
\end{Highlighting}
\end{Shaded}

Calculating the p-value is a two step process: first find the tail area,
and then multiply by two since the alternative hypothesis is two sided.

\begin{Shaded}
\begin{Highlighting}[]
\NormalTok{tail }\OtherTok{\textless{}{-}} \FunctionTok{pnorm}\NormalTok{(z, }\AttributeTok{lower.tail =} \ConstantTok{TRUE}\NormalTok{)}
\NormalTok{pval }\OtherTok{\textless{}{-}}\NormalTok{ tail }\SpecialCharTok{*} \DecValTok{2}
\end{Highlighting}
\end{Shaded}

Note that we used \texttt{lower.tail\ =\ TRUE} since the \(z\)-score we
obtained was negative. If we wanted to get the upper tail, which is
often the case when we have a positive \(z\)-score, we would use
\texttt{lower.tail\ =\ FALSE}.

\textbf{Exercise 8:} What is the conclusion of the hypothesis test? Make
sure to to provide an interpretation in context.

\textbf{Solution} If p-value \textless{} α, reject H0, data provide
evidence for HA If p-value \textgreater{} α, do not reject H0, data do
not provide evidence for HA If µ0 is inside the CI ⇒ fail to reject H0.
If µ0 is outside the CI ⇒ reject H0.

Since p-value=0.272 and is greater than 0.05 we fail to reject H0. The
data dose not provide convincing evidence that the average weight of
full term babies born in North Carolina is different than 7.5 pounds.

\subsection{More Practice}\label{more-practice}

\begin{enumerate}
\def\labelenumi{\arabic{enumi}.}
\tightlist
\item
  Using the original data set, \texttt{nc}, create a new data set called
  \texttt{nc\_sm} that only contains information from smoker mothers.
  Then calculate the average weight gained by these mothers.
\end{enumerate}

\begin{Shaded}
\begin{Highlighting}[]
\NormalTok{nc\_sm }\OtherTok{\textless{}{-}} \FunctionTok{subset}\NormalTok{(nc, nc}\SpecialCharTok{$}\NormalTok{habit }\SpecialCharTok{==} \StringTok{"smoker"}\NormalTok{)}
\end{Highlighting}
\end{Shaded}

Note that the variable `gained' contains missing values that will
produce a value of \texttt{NA} (i.e.~Not Available) in your
calculations. To remove the missing values (after you create the
\texttt{nc\_sm} data set), define a new variable called \texttt{gain}
using the following code :

\vspace{3mm}

\begin{Shaded}
\begin{Highlighting}[]
\NormalTok{gain }\OtherTok{\textless{}{-}}\NormalTok{ nc\_sm}\SpecialCharTok{$}\NormalTok{gained[}\FunctionTok{is.na}\NormalTok{(nc\_sm}\SpecialCharTok{$}\NormalTok{gained)}\SpecialCharTok{==}\NormalTok{F]}
\end{Highlighting}
\end{Shaded}

Use the \texttt{gain} variable you just created to perform all
calculations.

\begin{enumerate}
\def\labelenumi{\arabic{enumi}.}
\setcounter{enumi}{1}
\tightlist
\item
  Check if the assumptions and conditions necessary for inference on a
  mean are satisfied for this sample of smoker mothers.
\end{enumerate}

\textbf{Solution} We need to check that the data is a random sample,
independence, and the sampling distribution is approximately normal.

\begin{verbatim}
1. The given data was provided with a statement indicates that the sample is random, hence, the observations are independent.
2. We can see from the calculations below that the weight gained is within 10% of the mean. 
3. Since the sample size, 126 observations, is greater than 30, the Central Limit Theorem states the sampling distribution will be approximately normal regardless of the population shape. 
\end{verbatim}

\begin{Shaded}
\begin{Highlighting}[]
\NormalTok{nc\_sm }\OtherTok{\textless{}{-}} \FunctionTok{subset}\NormalTok{(nc, nc}\SpecialCharTok{$}\NormalTok{habit }\SpecialCharTok{==} \StringTok{"smoker"}\NormalTok{)}
\NormalTok{gain }\OtherTok{\textless{}{-}}\NormalTok{ nc\_sm}\SpecialCharTok{$}\NormalTok{gained[}\FunctionTok{is.na}\NormalTok{(nc\_sm}\SpecialCharTok{$}\NormalTok{gained)}\SpecialCharTok{==}\NormalTok{F]}
\NormalTok{xbar\_sm }\OtherTok{\textless{}{-}} \FunctionTok{mean}\NormalTok{(gain)}
\NormalTok{s\_sm }\OtherTok{\textless{}{-}} \FunctionTok{sd}\NormalTok{(gain)}
\NormalTok{n\_sm }\OtherTok{\textless{}{-}} \FunctionTok{length}\NormalTok{(gain)}
\NormalTok{boot\_mean\_sm }\OtherTok{\textless{}{-}} \FunctionTok{replicate}\NormalTok{(}\DecValTok{5000}\NormalTok{, }\FunctionTok{mean}\NormalTok{(}\FunctionTok{sample}\NormalTok{(gain, }\AttributeTok{replace =} \ConstantTok{TRUE}\NormalTok{)))}

\NormalTok{xbar\_sm}
\end{Highlighting}
\end{Shaded}

\begin{verbatim}
## [1] 31.92623
\end{verbatim}

\begin{Shaded}
\begin{Highlighting}[]
\FunctionTok{print}\NormalTok{(boot\_mean\_sm[}\DecValTok{1}\SpecialCharTok{:}\DecValTok{10}\NormalTok{])}
\end{Highlighting}
\end{Shaded}

\begin{verbatim}
##  [1] 31.98361 32.40984 30.55738 33.61475 31.42623 30.42623 29.37705 29.57377
##  [9] 31.50820 31.96721
\end{verbatim}

\begin{Shaded}
\begin{Highlighting}[]
\FunctionTok{hist}\NormalTok{(gain, }\AttributeTok{probability =} \ConstantTok{TRUE}\NormalTok{)}
\NormalTok{x\_sm }\OtherTok{=} \FunctionTok{seq}\NormalTok{(}\DecValTok{0}\NormalTok{, }\DecValTok{76}\NormalTok{, }\AttributeTok{by=}\NormalTok{.}\DecValTok{01}\NormalTok{)}
\NormalTok{y\_sm }\OtherTok{=} \FunctionTok{dnorm}\NormalTok{(}\AttributeTok{x =}\NormalTok{ x\_sm, }\AttributeTok{mean =} \FunctionTok{mean}\NormalTok{(gain), }\AttributeTok{sd =} \FunctionTok{sd}\NormalTok{(gain))}
\FunctionTok{lines}\NormalTok{(}\AttributeTok{x =}\NormalTok{ x\_sm, }\AttributeTok{y =}\NormalTok{ y\_sm, }\AttributeTok{col =} \StringTok{"blue"}\NormalTok{)}
\end{Highlighting}
\end{Shaded}

\begin{center}\includegraphics{Lab7-HW5_files/figure-latex/unnamed-chunk-17-1} \end{center}

\begin{Shaded}
\begin{Highlighting}[]
\FunctionTok{qqnorm}\NormalTok{(gain)}
\FunctionTok{qqline}\NormalTok{(gain)}
\end{Highlighting}
\end{Shaded}

\begin{center}\includegraphics{Lab7-HW5_files/figure-latex/unnamed-chunk-17-2} \end{center}

\begin{enumerate}
\def\labelenumi{\arabic{enumi}.}
\setcounter{enumi}{2}
\tightlist
\item
  Construct a 95\% confidence interval for the average weight gained by
  smoker mothers and provide an interpretation for this interval in
  context.
\end{enumerate}

\textbf{Solution} Since the confidence interval is always of the form

\begin{center}
$ point~estimate + ME$
\end{center}

Where the margin of error is always composed of the product of two
components, a critical value (t⋆ or z⋆) and a standard error.

\begin{Shaded}
\begin{Highlighting}[]
\NormalTok{se\_smoker }\OtherTok{\textless{}{-}}\NormalTok{ s\_sm }\SpecialCharTok{/} \FunctionTok{sqrt}\NormalTok{(n\_sm)}
\NormalTok{me\_smoker }\OtherTok{\textless{}{-}}\NormalTok{ zstar }\SpecialCharTok{*}\NormalTok{ se\_smoker}

\CommentTok{\# The bounds of the confidence interval can be calculated as}
\NormalTok{xbar\_sm }\SpecialCharTok{{-}}\NormalTok{ me\_smoker}
\end{Highlighting}
\end{Shaded}

\begin{verbatim}
## [1] 29.14897
\end{verbatim}

\begin{Shaded}
\begin{Highlighting}[]
\NormalTok{xbar\_sm }\SpecialCharTok{+}\NormalTok{ me\_smoker}
\end{Highlighting}
\end{Shaded}

\begin{verbatim}
## [1] 34.70349
\end{verbatim}

\begin{enumerate}
\def\labelenumi{\arabic{enumi}.}
\setcounter{enumi}{3}
\tightlist
\item
  Weight gain of roughly 30 pounds is considered normal. Conduct a
  hypothesis test to evaluate if the average weight gain of smoker
  mothers is different than 30 pounds. Make sure to interpret your
  conclusion in context of the question.
\end{enumerate}

\textbf{Solution} The Null Hypothesis (\(H_{0}\)): The average weight
gain of smoker mothers is equal to 30 pounds. Alternative Hypothesis
(\(H_{a}\)): The average weight gain of smoker mothers is not equal to
30 pounds. To conduct the hypothesis test we need to first calculate a
test statistic, a \(z\)-score, using the formula

\[ z = \frac{\bar{x} - \mu}{SE} \]

\begin{Shaded}
\begin{Highlighting}[]
\CommentTok{\# Calculate the following:}
\CommentTok{\# mu zero}
\CommentTok{\# number of smokers}
\CommentTok{\# standard deviation smokers}

\CommentTok{\# xbar\_sm \textless{}{-} mean(gain)   \#\# calculated above \^{}\^{}}
\NormalTok{mu\_sm }\OtherTok{\textless{}{-}} \DecValTok{30}
\NormalTok{s\_sm }\OtherTok{\textless{}{-}} \FunctionTok{sd}\NormalTok{(gain)}
\NormalTok{n\_sm }\OtherTok{\textless{}{-}} \FunctionTok{length}\NormalTok{(gain)}

\NormalTok{se\_sm }\OtherTok{\textless{}{-}}\NormalTok{ s\_sm }\SpecialCharTok{/} \FunctionTok{sqrt}\NormalTok{(n\_sm)}
\NormalTok{se\_sm}
\end{Highlighting}
\end{Shaded}

\begin{verbatim}
## [1] 1.416997
\end{verbatim}

\begin{Shaded}
\begin{Highlighting}[]
\NormalTok{z\_sm }\OtherTok{\textless{}{-}}\NormalTok{ (xbar\_sm }\SpecialCharTok{{-}}\NormalTok{ mu\_sm) }\SpecialCharTok{/}\NormalTok{ se\_sm}
\NormalTok{z\_sm}
\end{Highlighting}
\end{Shaded}

\begin{verbatim}
## [1] 1.359374
\end{verbatim}

Now we calculate a p-value: Since the alternative hypothesis is NOT two
sided, we only include one tail in the calculation for the p-value.

\begin{Shaded}
\begin{Highlighting}[]
\FunctionTok{pnorm}\NormalTok{(z\_sm, }\AttributeTok{lower.tail =} \ConstantTok{FALSE}\NormalTok{)}
\end{Highlighting}
\end{Shaded}

\begin{verbatim}
## [1] 0.087014
\end{verbatim}

\begin{Shaded}
\begin{Highlighting}[]
\CommentTok{\# tail \textless{}{-}pnorm(z\_sm, lower.tail = FALSE)}
\CommentTok{\# pval \textless{}{-} 2 * tail}
\CommentTok{\# pval}
\end{Highlighting}
\end{Shaded}

Since p-value=0.087014 is greater than 0.05 we fail to reject the Null
Hypothesis. The data dose not provide convincing evidence that the
average weight gain of mothers in North Carolina who smoke during
pregnancy is different than 30 pounds.

\begin{enumerate}
\def\labelenumi{\arabic{enumi}.}
\setcounter{enumi}{4}
\tightlist
\item
  Do the results of the confidence interval and hypothesis test agree
  with each other? Is this what you would expect or not? Explain.
\end{enumerate}

\textbf{Solution} Yes the results of the two tail hypothesis test and
the corresponding confidence interval will always agree; this is
expected since they are two different perspectives on the same
statistical question.

\subsubsection{Inference for
proportions}\label{inference-for-proportions}

Follow similar steps for the following exercise on inference for
proportions i.e.~check assumptions, show your calculations, state the
hypothesis, and interpret your results in context.

\begin{enumerate}
\def\labelenumi{\arabic{enumi}.}
\setcounter{enumi}{5}
\tightlist
\item
  Estimate the population proportion of premature births \(p\) using a
  95\% confidence interval.
\end{enumerate}

\textbf{Solution} Since the confidence interval is always of the form

\begin{center}
$ point~estimate + ME$
\end{center}

Where the margin of error is always composed of the product of a
critical value (t⋆ or z⋆) and the standard error. We will begin by
creating a value to store the number of premature births and construct a
95\% confidence interval \(z^\star\) to be

\begin{Shaded}
\begin{Highlighting}[]
\NormalTok{nc\_premi }\OtherTok{\textless{}{-}} \FunctionTok{subset}\NormalTok{(nc, nc}\SpecialCharTok{$}\NormalTok{premie }\SpecialCharTok{==} \StringTok{"premie"}\NormalTok{) }\CommentTok{\# North Carolina premature births}
\NormalTok{n\_premi }\OtherTok{\textless{}{-}} \FunctionTok{length}\NormalTok{(nc\_premi}\SpecialCharTok{$}\NormalTok{premie) }\CommentTok{\# Number of premature births}
\NormalTok{n\_total }\OtherTok{\textless{}{-}} \FunctionTok{length}\NormalTok{(nc}\SpecialCharTok{$}\NormalTok{premie) }\CommentTok{\# Total number in sample}

\NormalTok{p\_hat }\OtherTok{\textless{}{-}}\NormalTok{ n\_premi }\SpecialCharTok{/}\NormalTok{ n\_total }\CommentTok{\# the sample proportion}
\NormalTok{q\_hat }\OtherTok{\textless{}{-}} \DecValTok{1} \SpecialCharTok{{-}}\NormalTok{ p\_hat }\CommentTok{\# q = (1 {-} p)}

\CommentTok{\# Since we are calculating a 95\% CL, then calculate zstar}
\NormalTok{zstar }\OtherTok{\textless{}{-}} \FunctionTok{qnorm}\NormalTok{(}\FloatTok{0.975}\NormalTok{) }\DocumentationTok{\#\# equals 1.959964}
\end{Highlighting}
\end{Shaded}

The margin of error for a proportion is:

\$ MoE \$ = \$ (\{z\}\_\{\frac{\alpha }{2}\}
(\sqrt{\frac{\hat{p}\hat{q}}{n}}), where \hat{q}\text{ = 1 - }\hat{p} \$

\begin{Shaded}
\begin{Highlighting}[]
\CommentTok{\# premature \textless{}{-} length(nc$premie["premie"])}
\CommentTok{\# \# full\_term \textless{}{-} nc$premie["full term"]}
\CommentTok{\# n\_premature \textless{}{-} length(premature)}
\CommentTok{\# n\_total \textless{}{-} length(nc$premie)}
\CommentTok{\# n\_premature}
\CommentTok{\# n\_total}
\CommentTok{\# }
\CommentTok{\# }
\CommentTok{\# xbar\_premature \textless{}{-} mean(n\_premature)}
\CommentTok{\# }
\CommentTok{\# }
\CommentTok{\# \# premature \textless{}{-} nc$premie["premie"]}
\CommentTok{\# xbar\_premature \textless{}{-} mean(premature)}
\CommentTok{\# s\_premature \textless{}{-}sd(premature)}
\CommentTok{\# n\_premature \textless{}{-} length(premature)}
\CommentTok{\# }
\CommentTok{\# se\_premature \textless{}{-} s\_premature / sqrt(n\_premature)}
\CommentTok{\# me\_premature \textless{}{-} zstar * se\_premature}
\CommentTok{\# }
\CommentTok{\# \# The bounds of the confidence interval can be calculated as}
\CommentTok{\# xbar\_premature {-} me\_premature}
\CommentTok{\# xbar\_premature + me\_premature}
\end{Highlighting}
\end{Shaded}

The bounds of the confidence interval can be calculated as:

\begin{Shaded}
\begin{Highlighting}[]
\CommentTok{\# The bounds of the confidence interval can be calculated as}
\CommentTok{\# xbar\_premature {-} me\_premature}
\CommentTok{\# xbar\_premature + me\_premature}
\end{Highlighting}
\end{Shaded}

\begin{enumerate}
\def\labelenumi{\arabic{enumi}.}
\setcounter{enumi}{6}
\item
  Calculate a 95\% bootstrap confidence interval for the proportion of
  premature births. How does it compare to the interval in question 6?
\item
  In the United States, the proportion of live births that are premature
  is approximately 10\%. Test the hypothesis that the proportion of
  premature births in is greater than 10\%.
\end{enumerate}

\end{document}
